\section{Introduction}
Universal approximation leaves open how to find the weights of a particular small network. Consider a network given twelve examples of addition. How closely do the returned weights fit those sums, and do the same weights still add correctly when both numbers are larger? We evaluate each solver on both questions. The unit of comparison is a \emph{task and architecture}: fixing addition as the target still leaves a choice between a linear map with two coefficients and a network with a trainable hidden layer. The former has a convex least-squares training loss. The latter can have several weight vectors that fit equally well and a much worse fit halfway between them.

\emph{Weight-space programming} (WSP) uses a specification or examples to assign numerical values directly to a network's weights. Our solver, GPT-6 Astra (Astra), receives the architecture, a random initial weight vector and labeled training examples. In \emph{one-shot WSP}, it returns numerical weights for every trainable coefficient in a single reply. In \emph{iterative WSP}, we start with those very weights, evaluate them on training and development examples, and invite up to three revisions. We keep a revision only if it improves the development loss, using the training loss to break a tie. Each procedure is run with examples only and with the generating formula supplied alongside the same examples.

Backpropagation (BP) computes a derivative; its use does not specify how to update the weights. Our first numerical baseline uses that derivative for full-batch gradient descent. Our second begins at the same gradient-descent endpoint and uses a limited-memory approximation to local curvature to refine it. We report both checkpoints. Additional comparisons use adaptive gradient steps and trust-region least squares. We evaluate the weights on the original training data and on disjoint tests inside and outside the original input range. For integer sequences both tests concern unseen later indices.

The benchmark covers four elementary functions, two arithmetic operations and four sequence or sorting tasks. For every included fixed-network task, architecture and training set we exhibit two weight vectors whose midpoint has higher loss than the average endpoint loss. This strict violation of the convexity inequality certifies that the empirical objective is nonconvex in its trainable parameters. It permits the comparison to include linear recurrent networks whose repeated trainable matrices make their losses nonconvex. We measure how often each method fits the observations more closely, how often that advantage carries beyond the observed range, and how the outcome changes with the number of parameters. The question concerns the quality of the returned weights at stated budgets. The tasks have known mathematical answers and allow direct checks of exactness.

We also ask Astra to choose the architecture from examples. Thirty further requests specify the observations, deployment domain and supported neural operations. Astra chooses the graph, its size and coefficients, with extrapolation accuracy as the primary objective and inference cost as the secondary objective. We evaluate every returned proposal and report its coefficient storage, activation memory and arithmetic cost.


\input{holistic_main}
